\scanpage{153}
\margin{152}
\mspage{7.}

Pour compléter l'étude des relations entre le petit site $\operatorname{Ouv}(X)$ et le site $(\mathrm{Sch}/X,\mathrm{Zar})$, il reste à démontrer que chaque faisceau sur $\operatorname{Ouv}(X)$ est l'image par $i_{X\inserted{*}}$ d'un faisceau de Zariski sur $(\mathrm{Sch})_{/X}$. Ceci montrera en particulier, en vertu de l'isomorphisme (2), que la cohomologie de $\operatorname{Ouv}(X)$ pourra \inserted{ê}tre calculé dans $(\mathrm{Sch}/X,\mathrm{Zar})$. Plus précisément, nous avons le résultat suivant:

\uline{Proposition} 3.- \uline{Pour chaque} $F\inserted{\in}\overset{\inserted{\sim}}{\operatorname{Ouv}(X)}$, \uline{le morphisme naturel}
\[
F\inserted{\longrightarrow}i_{X\inserted{*}}i_X^{\inserted{*}}F
\]
\uline{est un isomorphisme.}

\uline{Démonstration:} Il suffit de démontrer que
\[
F(U)\inserted{\longrightarrow}(i_{X\bullet}ii_X^{\inserted{*}}F)(U)=(i_X^{\inserted{*}}F)(i_X(U))
\]
est un isomorphisme pour tout $U\inserted{\in}\operatorname{Ouv}(X)$. Il suffit m\inserted{ê}me de démontrer que
\[
F(U)\inserted{\longrightarrow}(i_X^\bullet F)(i_X(U))
\]
est un isomorphisme pour tout $U$. Rappelons que
\[
(i_X^\bullet F)(Y)=\varinjlim_{I_{i_X}^{Y\,o}}F(V)
\]
où $I_{i_X}^{Y}$ est la catégorie des paires $\inserted{\{}V,Y\inserted{\xrightarrow{\varphi}}i_X(V)\inserted{\}}$ $(Y\inserted{\in}(\mathrm{Sch})_{/X})$ avec, comme morphismes entre deux telles paires, les $V\inserted{\xrightarrow{j}}V'$ tels que le diagramme
\diagram{\begin{tikzcd}[ampersand replacement=\&]
\&Y\arrow[dl,"\inserted{\varphi}"']\arrow[dr,"\inserted{\varphi'}"]\&\\
i_X(V)\arrow[rr,"i_X(j)"]\&\&i_X(V')
\end{tikzcd}}
soit \uline{commutatif} (dans $(\mathrm{Sch})_{/X}$ ). Cette catégorie est pseudo-filtrante à gauche et connexe, car $i_X$ commute aux limites projectives finies. Dans le cas qui nous intéresse, $Y=i_X(U)$, il est facile de voir qu'il existe \inserted{en outre} un objet final dans la catégorie $I_{i_X}^{Y\,o}$, à savoir $\inserted{\{}U,i_X(U)\inserted{\xrightarrow{\mathrm{Id}}}i_X(U)\inserted{\}}$, d'où le résultat.
